\documentclass[11pt]{article}
\usepackage[margin=1in]{geometry}
\usepackage{amsmath, amssymb}
\usepackage{tikz}

\title{Every Vertex-Transitive Graph is Regular, but the Converse is Not True}
\author{Srijan Majumder \\ Scholar Id : 26CS22R04 }
\date{\today}

\begin{document}
\maketitle

\section*{\underline{Proof}}

A graph is \emph{regular} if every vertex has the same number of neighbours. A graph is \emph{vertex-transitive} if, for any two vertices $u$ and $v$, some symmetry of the graph (an automorphism) can slide $u$ exactly onto $v$, keeping every edge intact.

\subsection*{Vertex-transitive implies regular}

Suppose $G$ is vertex-transitive, and take any two vertices $u$ and $v$. There is an automorphism $f$ with $f(u) = v$. Since $f$ preserves edges, it matches each neighbour of $u$ to a neighbour of $v$, one-to-one, with none left over. So $u$ and $v$ have the same number of neighbours. As $u$ and $v$ were arbitrary, every vertex has the same degree -- $G$ is regular.

\subsection*{Regular does not imply vertex-transitive}

Regular only checks one number: the degree. Vertex-transitive checks that the whole neighbourhood looks the same everywhere, which is a stronger requirement. So a graph can match on degree alone without being truly symmetric.

\section*{\underline{Example}}

\begin{figure}[h]
\centering
\begin{tikzpicture}[scale=1, vtx/.style={circle,draw,fill=white,inner sep=1.5pt,minimum size=4pt}]
\node[vtx] (A1) at (0,1.6) {};
\node[vtx] (A2) at (-1.4,-0.7) {};
\node[vtx] (A3) at (1.4,-0.7) {};
\draw (A1) -- (A2) -- (A3) -- (A1);
\begin{scope}[xshift=6cm]
\foreach \i in {0,...,5} {
  \node[vtx] (B\i) at (90+60*\i:1.6) {};
}
\draw (B0) -- (B1) -- (B2) -- (B3) -- (B4) -- (B5) -- (B0);
\end{scope}
\end{tikzpicture}
\caption{A triangle and a hexagon, with no edges between them.}
\end{figure}

Take a triangle and a hexagon, drawn side by side with no edges connecting them.
Every vertex, in the triangle or the hexagon, has exactly 2 neighbours -- so the graph is regular.
But no symmetry can send a triangle vertex to a hexagon vertex: a symmetry must send a 3-cycle to a 3-cycle, and a hexagon vertex sits on a 6-cycle, not a 3-cycle. So the graph is not vertex-transitive.
This one picture is enough: regular, but not vertex-transitive.

\end{document}
